\begin{lemma}[Gradient annulus on a small ball]
Let $p\in\mathbb{R}$, let $F$ be a datum of Lewis--Tolksdorf type for $p$ as in
\noderef{FirstOrderEstimateInput}, with exponent $\kappa\in(0,1)$ and constant $C_0>0$, and
let $\delta>0$.  Then there is $\tau$ with $0<\tau<1/4$, depending only on $\delta$ and on the
data $\kappa,C_0$ (hence only on $p$ and $\delta$), with the following property.

Let $v\colon\mathrm{Plane}\to\mathbb{R}$ be $p$-harmonic on $B_1$ with weak gradient $g$ (in
the sense of \noderef{IsPHarm}), suppose $\lvert v\rvert\le1$ on $B_1$, and suppose
$\lvert Dv(0)\rvert>\delta$, where $Dv$ denotes the classical gradient.  Then $v$ is of class
$C^{1}$ on $B_{2\tau}$ and, for every $x\in B_{2\tau}$,
\[
  \frac{\delta}{2}\le\lvert Dv(x)\rvert\le C_0
  \qquad\text{and}\qquad
  \lvert Dv(x)-Dv(0)\rvert\le\frac{\delta}{2}.
\]

The third bound is the displacement estimate behind the first two, kept in the statement
because it is what a later Schauder bootstrap needs: it says that $Dv\bigl(B_{2\tau}\bigr)$ is
contained in the closed ball of radius $\delta/2$ about $Dv(0)$, which is \emph{convex} and
avoids the origin since $\lvert Dv(0)\rvert>\delta$.  The annulus
$\{\delta/2\le\lvert\xi\rvert\le C_0\}$ cut out by the first two bounds is not convex, so a
mean value inequality along a segment joining two values of $Dv$ cannot be run from those
alone.
\end{lemma}
\begin{proof}
\emph{Choice of $\tau$.}  Put
\[
  \tau:=\frac{1}{2}\min\Bigl\{\frac{1}{4},\;
    \Bigl(\frac{\delta}{2C_0}\Bigr)^{1/\kappa}\Bigr\} .
\]
Then $\tau>0$, because $\kappa>0$, $C_0>0$ and $\delta>0$ make both entries of the minimum
positive, and $\tau\le1/8<1/4$.  Moreover $2\tau\le(\delta/(2C_0))^{1/\kappa}$, so by
monotonicity of $t\mapsto t^{\kappa}$ on $[0,\infty)$ (clause (3) of the background proposition
on elementary properties of real powers) and the identity
$\bigl(s^{1/\kappa}\bigr)^{\kappa}=s$ for $s\ge0$ (clause (2) of that proposition),
\begin{equation}\label{eq:ga-tau}
  C_0\,(2\tau)^{\kappa}\le C_0\cdot\frac{\delta}{2C_0}=\frac{\delta}{2}.
\end{equation}
This is the paper's choice $C_0(2\tau)^{\kappa}\le\delta/2$.  The number $\tau$ was fixed from
$\delta,\kappa,C_0$ alone, before any function $v$ was given.

\emph{Verification.}  Let $v$ and $g$ be as in the statement.  Since $2\tau\le1/4<3/4$, we have
$B_{2\tau}\subseteq B_{3/4}$.  Applying the estimate carried by $F$ (clause \texttt{est} of
\noderef{FirstOrderEstimateInput}) to $v$, $g$ and the hypothesis $\lvert v\rvert\le1$ on
$B_1$, we obtain:
\begin{enumerate}
\item[(i)] $v$ is of class $C^{1}$ on $B_{3/4}$; by restriction (clause (4) of the background
proposition on locality of $C^{1}$ regularity) $v$ is of class $C^{1}$ on $B_{2\tau}$, which is
the first assertion;
\item[(ii)] $\lvert Dv(x)\rvert\le C_0$ for every $x\in B_{3/4}$, hence in particular for every
$x\in B_{2\tau}$, which is the upper bound;
\item[(iii)] $\lvert Dv(x)-Dv(y)\rvert\le C_0\lvert x-y\rvert^{\kappa}$ for all
$x,y\in B_{3/4}$.
\end{enumerate}

Now fix $x\in B_{2\tau}$.  Both $x$ and $0$ lie in $B_{3/4}$ (note $0\in B_{2\tau}$ because
$\tau>0$), so (iii) applies with $y=0$:
\[
  \lvert Dv(x)-Dv(0)\rvert\le C_0\,\lvert x\rvert^{\kappa}\le C_0\,(2\tau)^{\kappa}
  \le\frac{\delta}{2},
\]
where the middle inequality uses $\lvert x\rvert<2\tau$ and the monotonicity of
$t\mapsto t^{\kappa}$, and the last is \eqref{eq:ga-tau}.  By the triangle inequality in the
form $\lvert Dv(0)\rvert\le\lvert Dv(x)\rvert+\lvert Dv(x)-Dv(0)\rvert$ and the hypothesis
$\lvert Dv(0)\rvert>\delta$,
\[
  \lvert Dv(x)\rvert\ge\lvert Dv(0)\rvert-\lvert Dv(x)-Dv(0)\rvert
  >\delta-\frac{\delta}{2}=\frac{\delta}{2},
\]
so in particular $\lvert Dv(x)\rvert\ge\delta/2$.  Together with (ii) this is the asserted
two-sided bound.  The displayed estimate
$\lvert Dv(x)-Dv(0)\rvert\le C_0\lvert x\rvert^{\kappa}\le C_0(2\tau)^{\kappa}\le\delta/2$ just
used is itself the third asserted bound, so all three conclusions are established for the
arbitrary $x\in B_{2\tau}$ fixed above, and the proof is complete.
\end{proof}
